\ifdefined\gkver\else\def\gkver{student}\fi
\documentclass[\gkver]{gaokaozhenti}
\begin{document}

\chapter{2007年宁夏理综}
%% paperType: 理综物理部分

\begin{enumerate}
\item
%% number: 14
%% paperName: 2007年宁夏理综第 14 题 6 分
%% typeId: 1
%% type: 单选题
%% chapter: 曲线运动
%% point: 天体运动
%% method: 无
%% score: 6
%% degree: 650
%% duplicateId: 0
%% body:
天文学家发现了某恒星有一颗行星在圆形轨道上绕其运动，并测出了行星的轨道半径和运行周期。由此可推算出\xzanswer{C}

\fourchoices[answer=C]
{行星的质量}
{行星的半径}
{恒星的质量}
{恒星的半径}

%% answer: C
%% memo:
\memoanswer{
本题原卷及材料中未提供详解，待补充。
}

\item
%% number: 15
%% paperName: 2007年宁夏理综第 15 题 6 分
%% typeId: 1
%% type: 单选题
%% chapter: 曲线运动
%% point: 曲线运动条件
%% method: 无
%% score: 6
%% degree: 650
%% duplicateId: 0
%% body:
下列说法正确的是\xzanswer{B}

\fourchoices[answer=B]
{行星的运动和地球上物体的运动遵循不同的规律}
{物体在转弯时一定受到力的作用}
{月球绕地球运动时受到地球的引力和向心力的作用}
{物体沿光滑斜面下滑时受到重力、斜面的支持力和下滑力的作用}

%% answer: B
%% memo:
\memoanswer{
本题原卷及材料中未提供详解，待补充。
}

\item
%% number: 16
%% paperName: 2007年宁夏理综第 16 题 6 分
%% typeId: 1
%% type: 单选题
%% chapter: 直线运动
%% point: 运动图象
%% method: 无
%% score: 6
%% degree: 650
%% duplicateId: 0
%% body:
甲乙两辆汽车在平直的公路上沿同一方向作直线运动，$t=0$ 时刻同时经过公路旁的同一个路标。在描述两车运动的 $v-t$ 图中（如图），直线 a、b 分别描述了甲乙两车在 $0\sim20\Us$ 的运动情况。关于两车之间的位置关系，下列说法正确的是\xzanswer{C}

\onepicture[w=5.5cm]{figs/16.png}

\fourchoices[answer=C]
{在 $0\sim10\Us$ 内两车逐渐靠近}
{在 $10\sim20\Us$ 内两车逐渐远离}
{在 $5\sim15\Us$ 内两车的位移相等}
{在 $t=10\Us$ 时两车在公路上相遇}

%% answer: C
%% memo:
\memoanswer{
由 $v-t$ 图象知，$0\sim10\Us$ 内，$v_{\text{乙}}>v_{\text{甲}}$，两车逐渐远离；$10\sim20\Us$ 内，$v_{\text{乙}}<v_{\text{甲}}$，两车逐渐靠近，故 A、B 均错。$v-t$ 图线与时间轴所围的面积表示位移，$5\sim15\Us$ 内，两图线与 $t$ 轴包围的面积相等，故两车的位移相等，C 对。在 $t=10\Us$ 时，两车的位移不相等，说明两车不相遇，故 D 错。
}

\item
%% number: 17
%% paperName: 2007年宁夏理综第 17 题 6 分
%% typeId: 2
%% type: 多选题
%% chapter: 交流电
%% point: 交流电
%% method: 无
%% score: 6
%% degree: 650
%% duplicateId: 0
%% body:
一正弦交流电的电压随时间变化的规律如图所示。由图可知\xzanswer{BD}

\onepicture[w=7.0cm]{figs/17.png}

\fourchoices[answer=BD]
{该交流电的电压瞬时值的表达式为 $u=100\sin(25t)\UV$}
{该交流电的频率为 $25\UHz$}
{该交流电的电压的有效值为 $100\sqrt{2}\UV$}
{若将该交流电压加在阻值 $R=100\UO$ 的电阻两端，则电阻消耗的功率是 $50\UW$}

%% answer: BD
%% memo:
\memoanswer{
本题原卷及材料中未提供详解，待补充。
}

\item
%% number: 18
%% paperName: 2007年宁夏理综第 18 题 6 分
%% typeId: 1
%% type: 单选题
%% chapter: 电场
%% point: 带电粒子的直线运动
%% method: 整体与隔离
%% score: 6
%% degree: 450
%% duplicateId: 0
%% body:
两个质量相同的小球用不可伸长的细线连结，置于场强为 $E$ 的匀强电场中，小球 1 和小球 2 均带正电，电量分别为 $q_{1}$ 和 $q_{2}$（$q_{1}>q_{2}$）。将细线拉直并使之与电场方向平行，如图所示。若将两小球同时从静止状态释放，则释放后细线中的张力 $T$ 为\xzanswer{A}（不计重力及两小球间的库仑力）

\onepicture[w=5.0cm]{figs/18.png}

\fourchoices[answer=A]
{$T=\frac{1}{2}(q_{1}-q_{2})E$}
{$T=(q_{1}-q_{2})E$}
{$T=\frac{1}{2}(q_{1}+q_{2})E$}
{$T=(q_{1}+q_{2})E$}

%% answer: A
%% memo:
\memoanswer{
本题原卷及材料中未提供详解，待补充。
}

\item
%% number: 19
%% paperName: 2007年宁夏理综第 19 题 6 分
%% typeId: 1
%% type: 单选题
%% chapter: 电路
%% point: 动态分析
%% method: 无
%% score: 6
%% degree: 650
%% duplicateId: 0
%% body:
在如图所示的电路中，$E$ 为电源电动势，$r$ 为电源内阻，$R_{1}$ 和 $R_{3}$ 均为定值电阻，$R_{2}$ 为滑动变阻器。当 $R_{2}$ 的滑动触点在 a 端时合上开关 S，此时三个电表 A$_{1}$、A$_{2}$ 和 V 的示数分别为 $I_{1}$、$I_{2}$ 和 $U$。现将 $R_{2}$ 的滑动触点向 b 端移动，则三个电表示数的变化情况是\xzanswer{B}

\onepicture[w=5.0cm]{figs/19.png}

\fourchoices[answer=B]
{$I_{1}$ 增大，$I_{2}$ 不变，$U$ 增大}
{$I_{1}$ 减小，$I_{2}$ 增大，$U$ 减小}
{$I_{1}$ 增大，$I_{2}$ 减小，$U$ 增大}
{$I_{1}$ 减小，$I_{2}$ 不变，$U$ 减小}

%% answer: B
%% memo:
\memoanswer{
本题原卷及材料中未提供详解，待补充。
}

\item
%% number: 20
%% paperName: 2007年宁夏理综第 20 题 6 分
%% typeId: 1
%% type: 单选题
%% chapter: 电磁感应
%% point: 楞次定律
%% method: 无
%% score: 6
%% degree: 450
%% duplicateId: 0
%% body:
电阻 $R$、电容 $C$ 与一线圈连成闭合电路，条形磁铁静止于线圈的正上方，N 极朝下，如图所示。现使磁铁开始自由下落，在 N 极接近线圈上端的过程中，流过 $R$ 的电流方向和电容器极板的带电情况是\xzanswer{D}

\onepicture[w=4.2cm]{figs/20.png}

\fourchoices[answer=D]
{从 a 到 b，上极板带正电}
{从 a 到 b，下极板带正电}
{从 b 到 a，上极板带正电}
{从 b 到 a，下极板带正电}

%% answer: D
%% memo:
\memoanswer{
本题原卷及材料中未提供详解，待补充。
}

\item
%% number: 21
%% paperName: 2007年宁夏理综第 21 题 6 分
%% typeId: 1
%% type: 单选题
%% chapter: 电场
%% point: 等势面
%% method: 无
%% score: 6
%% degree: 450
%% duplicateId: 0
%% body:
匀强电场中的三点 A、B、C 是一个三角形的三个顶点，AB 的长度为 $1\Um$，D 为 AB 的中点，如图所示。已知电场线的方向平行于 $\triangle ABC$ 所在平面，A、B、C 三点的电势分别为 $14\UV$、$6\UV$ 和 $2\UV$。设场强大小为 $E$，一电量为 $1\times10^{-6}\UC$ 的正电荷从 D 点移到 C 点电场力所做的功为 $W$，则\xzanswer{A}

\onepicture[w=4.0cm]{figs/21.png}

\fourchoices[answer=A]
{$W=8\times10^{-6}\UJ$，$E>8\UVm$}
{$W=6\times10^{-6}\UJ$，$E>6\UVm$}
{$W=8\times10^{-6}\UJ$，$E\leq8\UVm$}
{$W=6\times10^{-6}\UJ$，$E\leq6\UVm$}

%% answer: A
%% memo:
\memoanswer{
本题原卷及材料中未提供详解，待补充。
}

\item
%% number: 22
%% paperName: 2007年宁夏理综第 22 题 0 分
%% typeId: 4
%% type: 实验题
%% chapter: 电路
%% point: 测电动势与内阻
%% method: 无
%% score: 0
%% degree: 650
%% duplicateId: 0
%% body:
22．实验题

\begin{enumerate}
	\item
	由绝缘介质隔开的两个同轴的金属圆筒构成圆柱形电容器，如图所示。试根据你学到的有关平行板电容器的知识，推测影响圆柱形电容器电容的因素有 \tkanswer[4cm]{$H$、$R_{1}$、$R_{2}$、$\varepsilon$（正对面积、板间距离、极板间的介质）}。

	\onepicture[h=5.5cm, align=right]{figs/22.png}

	\item
	利用伏安法测量干电池的电动势和内阻，现有的器材为：

	\begin{itemize}
		\item
		干电池：电动势约为 $1.5\UV$，符号如图\ref{2007宁夏理综22a}；
		\item
		电压表：量程 $1\UV$，内阻 $998.3\UO$，符号 V；
		\item
		电流表：量程 $1\UA$，符号 A；
		\item
		滑动变阻器：最大阻值 $99999.9\UO$，符号如图\ref{2007宁夏理综22b}；
		\item
		单刀单掷开关 1 个，符号如图\ref{2007宁夏理综22c}；
		\item
		导线若干。
	\end{itemize}

	\threepicture[labelA=2007宁夏理综22a, labelB=2007宁夏理综22b, labelC=2007宁夏理综22c, capA=干电池, capB=滑动变阻器, capC=单刀单掷开关, w=2.2cm, gap=1.6cm, align=right]
	{figs/22a.png}
	{figs/22b.png}
	{figs/22c.png}

	\begin{steps}
		\item
		设计测量电源电动势和内阻的电路并将它画在指定的方框内，要求在图中标出电压表、电流表的接线柱的正负。
		\item
		为了满足本实验要求并保证实验的精确度，电压表量程应扩大为原量程的 \tkanswer[1cm]{2} 倍，电阻箱的阻值应为 \tkanswer[1.5cm]{998.3} $\UO$。
	\end{steps}
\end{enumerate}

%% answer:
\jdanswer{
\begin{enumerate}
	\item
	$H$、$R_{1}$、$R_{2}$、$\varepsilon$（正对面积、板间距离、极板间的介质）。

	\item
	①电路图如图所示。

	\onepicture[w=6.0cm]{figs/22_an.png}

	②2；$998.3\UO$。
\end{enumerate}
}

%% memo:
\memoanswer{
本题原卷及材料中未提供详解，待补充。
}

\item
%% number: 23
%% paperName: 2007年宁夏理综第 23 题 0 分
%% typeId: 6
%% type: 计算题
%% chapter: 曲线运动
%% point: 平抛运动
%% method: 无
%% score: 0
%% degree: 450
%% duplicateId: 0
%% body:
倾斜雪道的长为 $25\Um$，顶端高为 $15\Um$，下端经过一小段圆弧过渡后与很长的水平雪道相接，如图所示。一滑雪运动员在倾斜雪道的顶端以水平速度 $v_{0}=8\Ums$ 飞出，在落到倾斜雪道上时，运动员靠改变姿势进行缓冲使自己只保留沿斜面的分速度而不弹起。除缓冲外运动员可视为质点，过渡轨道光滑，其长度可忽略。设滑雪板与雪道的动摩擦因数 $\mu=0.2$，求运动员在水平雪道上滑行的距离。（取 $g=10\Umsq$）

\onepicture[w=7.0cm, align=right]{figs/23.png}

%% answer:
\jdanswer{
$s_{2}=74.8\Um$
}

%% memo:
\memoanswer{
如图选坐标，斜面的方程为 $y=x\tan\theta=\frac{3}{4}x$（①）。

\onepicture[w=5.0cm]{figs/23_an.png}

运动员飞出后做平抛运动 $x=v_{0}t$（②），$y=\frac{1}{2}gt^{2}$（③）。联立①②③式，得飞行时间 $t=1.2\Us$，落点的 $x$ 坐标 $x_{1}=v_{0}t=9.6\Um$，落点离斜面顶端的距离 $s_{1}=\frac{x}{\cos\theta}=12\Um$，落点距地面的高度 $h_{1}=(L-s_{1})\sin\theta=7.8\Um$。接触斜面前的 $x$ 分速度 $v_{x}=8\Ums$，$y$ 分速度 $v_{y}=gt=12\Ums$，沿斜面的速度大小为 $v_{B}=v_{x}\cos\theta+v_{y}\sin\theta=13.6\Ums$。设运动员在水平雪道上运动的距离为 $s_{2}$，由功能关系得 $mgh_{1}+\frac{1}{2}mv_{B}^{2}=\mu mg\cos\theta(L-s_{1})+\mu mgs_{2}$，解得 $s_{2}=74.8\Um$。
}

\item
%% number: 24
%% paperName: 2007年宁夏理综第 24 题 0 分
%% typeId: 6
%% type: 计算题
%% chapter: 磁场
%% point: 带电粒子在磁场中的运动
%% method: 无
%% score: 0
%% degree: 450
%% duplicateId: 0
%% body:
在半径为 $R$ 的半圆形区域中有一匀强磁场，磁场的方向垂直于纸面，磁感应强度为 $B$。一质量为 $m$，带有电量 $q$ 的粒子以一定的速度沿垂直于半圆直径 AD 方向经 P 点（$AP=d$）射入磁场（不计重力影响）。

\begin{enumerate}
	\item
	如果粒子恰好从 A 点射出磁场，求入射粒子的速度。
	\item
	如果粒子经纸面内 Q 点从磁场中射出，出射方向与半圆在 Q 点切线方向的夹角为 $\varphi$（如图）。求入射粒子的速度。
\end{enumerate}

\onepicture[w=5.0cm, align=right]{figs/24.png}

%% answer:
\jdanswer{
\begin{enumerate}
	\item
	$v_{1}=\frac{qBd}{2m}$
	\item
	$v=\frac{qBd(2R-d)}{2m[R(1+\cos\varphi)-d]}$
\end{enumerate}
}

%% memo:
\memoanswer{
（1）由于粒子在 P 点垂直射入磁场，故圆弧轨道的圆心在 AP 上，AP 是直径。设入射粒子的速度为 $v_{1}$，由洛伦兹力的表达式和牛顿第二定律得 $m\frac{v_{1}^{2}}{d/2}=qv_{1}B$，解得 $v_{1}=\frac{qBd}{2m}$。

（2）设 $O'$ 是粒子在磁场中圆弧轨道的圆心，连接 $O'Q$，设 $O'Q=R'$。由几何关系得 $\angle OQO'=\varphi$，$OO'=R+R'-d$，由余弦定理得 $(OO')^{2}=R^{2}+R'^{2}-2RR'\cos\varphi$，解得 $R'=\frac{d(2R-d)}{2[R(1+\cos\varphi)-d]}$。设入射粒子的速度为 $v$，由 $m\frac{v^{2}}{R'}=qBv$，解出 $v=\frac{qBd(2R-d)}{2m[R(1+\cos\varphi)-d]}$。

\onepicture[w=6.0cm]{figs/24_an.png}
}

\item
%% number: 30
%% paperName: 2007年宁夏理综第 30 题 0 分
%% typeId: 6
%% type: 计算题
%% chapter: 运动和力
%% point: 动量守恒定律
%% method: 无
%% score: 0
%% degree: 650
%% duplicateId: 0
%% body:
30．物理选考题

\begin{enumerate}
	\item
	（A）【物理——选修 2—2】

	塔式起重机的结构如图所示，设机架重 $P=400\UkN$，悬臂长度为 $L=10\Um$，平衡块重 $W=200\UkN$，平衡块与中心线 $OO'$ 的距离可在 $1\Um$ 到 $6\Um$ 间变化，轨道 A、B 间的距离为 $4\Um$。

	（1）当平衡块离中心线 $1\Um$，右侧轨道对轮子的作用力 $f_{B}$ 是左侧轨道对轮子作用力 $f_{A}$ 的 2 倍，问机架重心离中心线的距离是多少？

	（2）当起重机挂钩在离中心线 $OO'$ $10\Um$ 处吊起重为 $G=100\UkN$ 的重物时，平衡块离 $OO'$ 的距离为 $6\Um$，问此时轨道 B 对轮子的作用力 $F_{B}$ 是多少？

	\onepicture[w=4.5cm, align=right]{figs/30a.png}

	\item
	（B）【物理——选修 3—3】

	如图所示，两个可导热的气缸竖直放置，它们的底部都由一细管连通（忽略细管的容积）。两气缸各有一个活塞，质量分别为 $m_{1}$ 和 $m_{2}$，活塞与气缸无摩擦。活塞的下方为理想气体，上方为真空。当气体处于平衡状态时，两活塞位于同一高度 $h$。（已知 $m_{1}=3m$，$m_{2}=2m$）

	（1）在两活塞上同时各放一质量为 $m$ 的物块，求气体再次达到平衡后两活塞的高度差（假定环境温度始终保持为 $T_{0}$）。

	（2）在达到上一问的终态后，环境温度由 $T_{0}$ 缓慢上升到 $T$，试问在这个过程中，气体对活塞做了多少功？气体是吸收还是放出了热量？（假定在气体状态变化过程中，两物块均不会碰到气缸顶部）。

	\onepicture[w=5.0cm, align=right]{figs/30b.png}

	\item
	（C）【物理——选修 3—4】

	图为沿 $x$ 轴向右传播的简谐横波在 $t=1.2\Us$ 时的波形，位于坐标原点处的观察者测到在 $4\Us$ 内有 10 个完整的波经过该点。

	（1）求该波的波幅、频率、周期和波速。

	（2）画出平衡位置在 $x$ 轴上 P 点处的质点在 $0\sim0.6\Us$ 内的振动图象。

	\onepicture[w=8.5cm, align=right]{figs/30c.png}

	\item
	（D）【物理——选修 3—5】

	在光滑的水平面上，质量为 $m_{1}$ 的小球 A 以速率 $v_{0}$ 向右运动。在小球的前方 O 点处有一质量为 $m_{2}$ 的小球 B 处于静止状态，如图所示。小球 A 与小球 B 发生正碰后小球 A、B 均向右运动。小球 B 被在 Q 点处的墙壁弹回后与小球 A 在 P 点相遇，$PQ=1.5PO$。假设小球间的碰撞及小球与墙壁之间的碰撞都是弹性的，求两小球质量之比 $\frac{m_{1}}{m_{2}}$。

	\onepicture[w=7.0cm, align=right]{figs/30d.png}
\end{enumerate}

%% answer:
\jdanswer{
\begin{enumerate}
	\item
	（1）$x=1.5\Um$；（2）$F_{B}=450\UkN$。
	\item
	（1）$x=\frac{5}{4}h$（即两活塞的高度差为 $\frac{5}{4}h$）；（2）$W=5mgh\left(\frac{T}{T_{0}}-1\right)$，在此过程中气体吸收热量。
	\item
	（1）$A=0.1\Um$，$f=2.5\UHz$，$T=0.4\Us$，$v=5\Ums$；（2）振动图象如图所示。
	\item
	$\frac{m_{1}}{m_{2}}=2$
\end{enumerate}
}

%% memo:
\memoanswer{
A、解：（1）空载时合力为零 $f_{A}+f_{B}=P+W=600\UkN$，已知 $f_{B}=2f_{A}$，求得 $f_{A}=200\UkN$，$f_{B}=400\UkN$。设机架重心在中心线右侧，离中心线的距离为 $x$，以 A 为转轴，力矩平衡 $f_{B}\times4=W\times(2-1)+P\times(2+x)$，求得 $x=1.5\Um$。

（2）以 A 为转轴，力矩平衡 $W\times(6-2)+F_{B}\times4=P\times(2+1.5)+G\times(10+2)$，求得 $F_{B}=450\UkN$。

B、解：（1）设左、右活塞的面积分别为 $A'$ 和 $A$，由于气体处于平衡状态，故两活塞对气体的压强相等，即 $\frac{3mg}{A'}=\frac{2mg}{A}$，由此得 $A'=\frac{3}{2}A$。在两个活塞上各加一质量为 $m$ 的物块后，右活塞降至气缸底部，所有气体都在左气缸中。在初态，气体的压强为 $\frac{2mg}{A}$，体积为 $\frac{5}{2}Ah$；在末态，气体压强为 $\frac{8mg}{3A}$，体积为 $\frac{3}{2}Ax$（$x$ 为左活塞的高度）。由玻意耳定律得 $\frac{2mg}{A}\cdot\frac{5}{2}Ah=\frac{8mg}{3A}\cdot\frac{3}{2}Ax$，解得 $x=\frac{5}{4}h$，即两活塞的高度差为 $\frac{5}{4}h$。

（2）当温度由 $T_{0}$ 上升至 $T$ 时，气体的压强始终为 $\frac{8mg}{3A}$，设 $x'$ 是温度达到 $T$ 时左活塞的高度，由盖·吕萨克定律得 $x'=\frac{T}{T_{0}}x=\frac{5Th}{4T_{0}}$，气体对活塞做的功为 $W=Fs=4mg\cdot\left(\frac{5Th}{4T_{0}}-\frac{5h}{4}\right)=5mgh\left(\frac{T}{T_{0}}-1\right)$，在此过程中气体吸收热量。

C、解：（1）$A=0.1\Um$，$f=\frac{n}{t}=2.5\UHz$，$T=\frac{1}{f}=0.4\Us$，$v=\lambda f=5\Ums$。

（2）振动图象如图所示。

\onepicture[w=6.5cm]{figs/30c_an.png}

D、解：从两小球碰撞后到它们再次相遇，小球 A 和 B 的速度大小保持不变，根据它们通过的路程，可知小球 B 和小球 A 在碰撞后的速度大小之比为 $4:1$。两球碰撞过程有 $m_{1}v_{0}=m_{1}v_{1}+m_{2}v_{2}$，$\frac{1}{2}m_{1}v_{0}^{2}=\frac{1}{2}m_{1}v_{1}^{2}+\frac{1}{2}m_{2}v_{2}^{2}$，解得 $\frac{m_{1}}{m_{2}}=2$。
}

\end{enumerate}

\end{document}
